\begin{textsample}{POS Dim 2 – ba – Score 26.00 – ba/ba\_em\_2.txt}  \label{ex:f2_pos_037}
2.
INTRODUÇÃO “ Vem, vamos embora que esperar não é saber, quem sabe faz a hora, não espera acontecer ”. ( Geraldo Vandré ) O debate em torno do Ensino Médio brasileiro, influenciado pelos movimentos educacionais globalizados, vem ganhando projeção na pauta governamental, internacional e nacional, nos últimos anos.
No Brasil, atualmente, a \textbf{política} curricular tem como marco o Projeto de Lei ( PL ) nº 6.840/13, que propunha, entre outras disposições, \textbf{instituir} “ a jornada em tempo integral no ensino médio, dispor sobre a organização dos \textbf{currículos} em áreas do conhecimento ” ( BRASIL, 2013 ).
A partir do referido PL, o Conselho Nacional de \textbf{Secretários} de Educação ( CONSED ), constituiu Grupos de Trabalho ( GT ) temáticos voltados para estudos dos temas centrais para a \textbf{reforma} do Ensino Médio brasileiro.
Os GTs foram integrados pelos/as secretários/as \textbf{estaduais} de educação e técnicos/as dos setores responsáveis pela coordenação do Ensino Médio das Secretarias de Educação, das 27 ( vinte e sete ) Unidades Federativas ( UF ) e do Distrito Federal ( DF ).
Em 2014, com a sanção da Lei nº 13.005/14, que aprovou o Plano Nacional de Educação ( PNE ), foram aprovadas estratégias para subsidiar no cumprimento da meta 3, que estabelece : > Universalizar, até 2016, o atendimento escolar para toda a população de 15 ( quinze ) a 17 ( dezessete ) anos e elevar, até o final do período de vigência deste PNE, a taxa líquida de matrículas no Ensino Médio para 85 \% ( oitenta e cinco por cento ). ( BRASIL, 2016 ) As estratégias elaboradas direcionaram as ações prioritárias \textbf{estaduais} para que a meta fosse cumprida, dentre elas, algumas estratégias foram \textbf{destinadas} à reorganização curricular do Ensino Médio, a saber : institucionalização de programa nacional de renovação do ensino médio, de incentivo a práticas pedagógicas com abordagens interdisciplinares integrando teoria e prática, por meio de \textbf{currículos} escolares que organizem, de maneira \textbf{flexível} e diversificada, conteúdos obrigatórios e \textbf{eletivos} articulados em dimensões como ciência, trabalho, linguagens, tecnologia, cultura e esporte ; e a definição de direitos e objetivos de aprendizagem e desenvolvimento para os ( as ) alunos ( as ) de ensino médio, que configurassem a elaboração de uma Base Nacional Comum Curricular ( BNCC ) do Ensino Médio.
Neste sentido, a primeira versão da BNCC, que integrou todas as etapas da Educação Básica - Educação Infantil, Ensino Fundamental e Ensino Médio, organizada em uma perspectiva de progressão das aprendizagens dos/as estudantes fora apresentada sob a forma de consulta pública à sociedade brasileira.
A segunda versão, ajustada após consulta pública, foi aperfeiçoada durante a realização dos 27 seminários \textbf{estaduais} da BNCC, organizados pelo Ministério da Educação ( MEC ).
A versão final da BNCC, contemplando as três etapas da Educação Básica, seguiria para a apreciação do Conselho Nacional da Educação ( CNE ) e, caso aprovada, seria \textbf{homologada} pelo MEC.
Contudo, ao final do segundo \textbf{semestre} de 2016, com a mudança do governo da presidente Dilma Rousseff para o governo de Michel Temer, a proposta do PL nº 6.840/13, com alterações substanciais ao texto original se configura na Medida Provisória ( MP ) nº 746/16 e, em fevereiro de 2017, é aprovada se tornando na Lei nº 13.415/2017, que alterou a Lei de \textbf{Diretrizes} e Base da Educação ( LDB ), sobretudo nos artigos voltados para a organização dos \textbf{currículos} do Ensino Médio e por conseguinte se fez necessário a elaboração de uma nova versão da BNCC da etapa do Ensino Médio.
A nova versão da BNCC foi produzida, \textbf{alinhada} a Lei nº 13.415/17, e foi encaminhada ao Conselho Nacional de Educação ( CNE ).
O CNE organizou audiências públicas nas cinco regiões da federação para subsidiar a análise do documento e assim a BNCC do Ensino Médio, com as contribuições dos territórios, foi aprovada pelo CNE e \textbf{homologada} pelo MEC, em dezembro de 2018.
Além da BNCC, outros normativos e documentos foram produzidos no âmbito do CNE e MEC para orientar e apoiar os Sistemas, \textbf{Redes} e instituições de ensino no processo de implementação da \textbf{reforma} do Ensino Médio brasileiro, tais como : - \textbf{Portaria} MEC nº 649, de 10 de julho de 2018, que \textbf{institui} o Programa de Apoio ao Novo Ensino Médio. - \textbf{Portaria} MEC nº 1.023, de 4 de outubro de 2018, que estabelece \textbf{diretrizes}, parâmetros e critérios para a realização de avaliação de impacto do Programa de Fomento às Escolas de Ensino Médio em Tempo Integral ( EMTI ) e seleção de novas unidades escolares para o Programa. - \textbf{Portaria} MEC nº 1.024, de 4 de outubro de 2018, que define as \textbf{diretrizes} do apoio financeiro por meio do Programa Dinheiro Direto na Escola ( PDDE ) às unidades escolares pertencentes às Secretarias participantes do Programa de Apoio ao Novo Ensino Médio. - Resolução CNE/CEB nº 3, de 21 de novembro de 2018, que \textbf{atualiza} as \textbf{Diretrizes} Curriculares Nacionais do Ensino Médio. - \textbf{Portaria} MEC nº 1.432, de 28 de dezembro de 2018, que estabelece os referenciais para a elaboração dos \textbf{Itinerários} \textbf{Formativos}.
Não obstante à \textbf{legislação}, ressalta -se as intencionalidades que estruturam as propostas de \textbf{reformas} educacionais globalizadas que operam dentro da lógica da manutenção do modo de produção capitalista.
O panorama das \textbf{reformas} realizadas na maioria dos países da América Latina se iniciou nos anos 90, em contextos das \textbf{reformas} neoliberais, promulgadas pelos Estados nacionais em distintas partes do mundo (... ) por meio da gestão e da descentralização da educação, do financiamento e da estruturação dos sistemas educacionais, da avaliação e do \textbf{currículo} ( COSTA, 2011 ).
Os principais marcos indutores dessas \textbf{reformas} foram a Conferência Mundial de Educação para Todos, realizada em Jomtien, em 1990, que resultou na assinatura da Declaração Mundial sobre Educação para Todos e o Marco de Ação para a Satisfação das Necessidades Básicas de Aprendizagem, em 1993, ( MALANCHEN, 2016 ).
Além desses, outro documento de projeção internacional, a serviço do capital e crucial no movimento reformista da educação no Brasil, foi o Relatório da Comissão Internacional sobre Educação para o século XXI, elaborado pela Organização das Nações Unidas para a Educação, a Ciência e a Cultura ( UNESCO ).
No Brasil, o documento foi publicado em 1998 com o título “ Educação : um tesouro a descobrir ”, fruto de uma parceria entre a UNESCO e o Ministério da Educação ( MEC ).
Frente a esse cenário, propor \textbf{políticas} públicas educacionais para os/as trabalhadores/as e seus/suas filhos e filhas, que promovam a superação da lógica do capital é bastante desafiador, pois muitos elementos que estruturam a educação nacional estão a serviço da manutenção da sociedade capitalista de consumo e da divisão de classes sociais, a exemplo da difusão crescente de teorias que reforçam a ideia da escola como aparelho ideológico do estado, bem como o financiamento da educação no país é praticado por organismos internacionais.
Nesse contexto, a Bahia, ao longo de sua história, marcada por lutas de resistência e movimentos de libertação do seu povo e da sua ancestralidade, na atualidade, é arriscado qualquer passo que leve a retrocessos.
Sendo assim, deve ser um compromisso político deste Estado propor \textbf{políticas} que promovam a reparação e \textbf{equidade} social, inclusive as \textbf{políticas} educacionais como a \textbf{política} curricular que está sendo tratada neste documento.
Partindo da premissa que o papel social da educação é emancipar pessoas, se entende a educação escolar como > um processo ao qual compete oportunizar a apropriação do conhecimento historicamente sistematizado – o enriquecimento do universo de significações-, tendo em vista a \textbf{elevação} para além das significações mais imediatas e aparentes disponibilizadas pelas dimensões meramente empíricas dos fenômenos. ( Martins, 2013 ).
É importante ressaltar que apenas a \textbf{reforma} na \textbf{política} curricular para o Ensino Médio não poderá promover transformações profundas no cenário educacional e social dos/as estudantes desta etapa de ensino.
As transformações mais significativas só serão possíveis quando houver a integralização entre \textbf{políticas} públicas voltadas para a \textbf{garantia} dos direitos de todos/as, incluindo as \textbf{políticas} educacionais e políticas para as juventudes, sobretudo, para as juventudes negras, que estão em um contexto de marginalização histórica no país.
Nesse excerto, o que a Bahia visa, em uma perspectiva sócio-histórica, é apresentar um Referencial Curricular que se fundamenta em uma educação que promova o desenvolvimento do gênero humano, por meio da apropriação dos conhecimentos, historicamente construídos pela humanidade, sem deixar de promover a valorização e preservação de saberes e conhecimentos da cultura e tradições dos agrupamentos humanos em que vivem, nos 27 ( vinte e sete ) Territórios de Identidade do estado.
Este Volume 2 do Documento Curricular Referencial da Bahia – etapa do Ensino Médio retoma e aprofunda categorias já tratadas no Volume 1, como : os pressupostos legais, os temas integradores, a territorialidade, o compromisso com a progressão das aprendizagens dos/das estudantes e o Projeto de Vida considerados na transição entre o Ensino Fundamental - Anos Finais e o Ensino Médio.
Assim sendo, se faz necessário que os/as \textbf{profissionais} da educação se apropriem do Volume 1 do DCRB, sobretudo dos Anos Finais do Ensino Fundamental, para terem conhecimento do que foi considerado para a aprendizagem dos/as estudantes.
Nessa perspectiva, de progressão de aprendizagens, \textbf{preconizada} pelas primeiras versões da BNCC, antes do advento da Lei nº 13.415/17, é que os Volumes 1, 2 e 3 do DCRB buscam retomar essa premissa, estabelecendo nexos de integração entre os três volumes.
Destarte, convidamos os/as \textbf{profissionais} da educação, estudantes, familiares e demais segmentos da sociedade civil a se apropriarem deste Documento Curricular Referencial da Bahia ( DCRB ), etapa do Ensino Médio, para conhecimento das suas intencionalidades políticas e pedagógicas propostas às escolas baianas, como diretriz curricular para a formação integral dos/as estudantes do Ensino Médio da Bahia.

% matched lemmas: alinhar, atualizar, currículo, destinar, diretriz, eletivo, elevação, equidade, estadual, flexível, formativo, garantia, homologar, instituir, itinerário, legislação, política, portaria, preconizar, profissional, rede, reforma, secretário, semestre
\end{textsample}
